\section{Metodologia}

A presente etapa do projeto concentrou-se na consolidação teórica e na
implementação de uma prova de conceito computacional voltada à aplicação
de algoritmos quânticos variacionais na simulação de sistemas moleculares.
Dessa forma, a metodologia adotada foi estruturada de modo a integrar
capacitação conceitual, domínio de ferramentas computacionais e
implementação prática do algoritmo Variational Quantum Eigensolver (VQE),
permitindo a compreensão completa do fluxo híbrido quântico-clássico
característico dos algoritmos da era NISQ.

\subsection{Consolidação teórica e capacitação técnica}

Inicialmente, realizou-se um processo sistemático de estudo dos fundamentos
da computação quântica, algoritmos variacionais e métodos de simulação
molecular. Essa etapa constituiu parte central das atividades desenvolvidas,
uma vez que o objetivo principal desta fase do projeto consistiu na
consolidação conceitual necessária para futuras aplicações em problemas
de interesse na área de materiais energéticos.

A consolidação dos fundamentos de computação quântica foi realizada por meio 
da disciplina ``Computação Quântica'', ofertada na Unidade Acadêmica de Belo 
Jardim (UFRPE) e ministrada pelo orientador do projeto, Prof. Dr. Fábio 
Magalhães de Novaes Santos. 
Os fundamentos foram consolidados por meio de formação formal na disciplina 
desde os fundamentos matemáticos da mecânica quântica, incluindo 
notação de Dirac e representação vetorial de estados 
quânticos, até algoritmos quânticos fundamentais, como os algoritmos de 
Grover e Shor, além de tópicos avançados relacionados à simulação quântica e 
caminhadas quânticas, por exemplo. Essa formação ocorreu simultaneamente ao 
desenvolvimento da iniciação científica, contribuindo diretamente para a 
compreensão dos modelos computacionais empregados neste trabalho.

Foram utilizados como referências principais textos clássicos da área de
computação quântica \cite{nielsen2000quantum,hidary2019quantum}, bem como
artigos de revisão sobre algoritmos variacionais e computação quântica
aplicada à química molecular. Paralelamente, foi realizado o acompanhamento
do curso educacional \textit{Quantum Chemistry with VQE}, disponibilizado
pela plataforma IBM Quantum Learning \cite{ibm_vqe_course}, o qual apresenta
de forma estruturada as etapas necessárias para implementação do algoritmo
VQE em problemas de estrutura eletrônica. 
Essa etapa permitiu o domínio conceitual dos elementos fundamentais do
pipeline variacional, incluindo construção do Hamiltoniano molecular,
definição de estados variacionais, estratégias de medição e métodos de
otimização clássica.

\subsection{Ambiente computacional e ferramentas utilizadas}

Conforme previsto no plano de trabalho inicial, foram consideradas as 
plataformas Ket, QuBox e Qiskit para a realização das simulações quânticas. 
Entretanto, durante a fase de implementação, observou-se que o ecossistema 
desenvolvido pela IBM Quantum, baseado no framework Qiskit, apresenta 
integração completa entre simulação local, execução em hardware quântico 
real e ferramentas de química quântica variacional.
Dessa forma, optou-se pela adoção prioritária do Qiskit como ambiente 
unificado de desenvolvimento, permitindo a execução de todas as etapas 
necessárias ao projeto sem perda metodológica ou experimental. Essa escolha 
reduziu a complexidade do fluxo computacional e favoreceu uma maior 
padronização dos experimentos realizados.

As simulações foram executadas exclusivamente em simuladores quânticos
clássicos disponibilizados pelo Qiskit Aer, possibilitando a análise
do comportamento ideal do algoritmo na ausência de ruídos experimentais.
Essa escolha permitiu isolar aspectos conceituais e metodológicos do
algoritmo, priorizando a compreensão do seu funcionamento interno.

\subsection{Formulação do problema molecular}

Como sistemas modelo para implementação do algoritmo foram selecionadas
as moléculas de hidrogênio (H$_2$) e hidreto de lítio (LiH), amplamente
empregadas na literatura como benchmarks em simulações quânticas devido
à sua baixa complexidade computacional e estrutura eletrônica bem
caracterizada.

\subsubsection{Definição geométrica e base atômica}

As geometrias moleculares foram definidas explicitamente por meio
das coordenadas cartesianas dos núcleos atômicos. No caso da
molécula de hidrogênio, os átomos foram posicionados
simetricamente em relação à origem, separados por uma distância
internuclear de 0,735 Å, valor próximo ao comprimento de ligação
de equilíbrio experimental.
Para ambas as moléculas foi adotada a base atômica STO-6G,
base mínima amplamente utilizada em estudos introdutórios
por reduzir o custo computacional mantendo as características
essenciais da estrutura eletrônica.

\subsubsection{Cálculo eletrônico clássico e espaço ativo}

A estrutura eletrônica inicial foi obtida por meio do método
Hartree–Fock restrito (RHF), implementado na biblioteca PySCF.
Em seguida, foi empregado o método CASCI (Complete Active Space
Configuration Interaction), permitindo definir explicitamente
um espaço ativo contendo os orbitais moleculares mais relevantes
para a descrição da correlação eletrônica.

No caso da molécula de hidrogênio, foram considerados dois
orbitais ativos com dois elétrons no espaço ativo,
enquanto para o LiH foi adotado um espaço ativo ampliado,
compatível com sua maior complexidade eletrônica.
A seleção dos orbitais ativos foi realizada com base na
simetria molecular, garantindo consistência física na
descrição do sistema.

\subsubsection{Construção do Hamiltoniano eletrônico}

A partir do cálculo CASCI foram extraídos:

\begin{itemize}
\item os integrais de um elétron ($h_{pq}$);
\item os integrais de dois elétrons ($h_{pqrs}$);
\item a energia de núcleo efetiva ($E_{\text{core}}$).
\end{itemize}

Os termos de dois elétrons foram submetidos a uma decomposição
de Cholesky de baixo posto, procedimento que permite fatorar o
tensor eletrônico reduzindo a complexidade computacional da
construção do Hamiltoniano.
Essa etapa resulta em uma representação eficiente do Hamiltoniano
fermiónico, posteriormente mapeado para operadores de spin por
meio da transformação de Jordan–Wigner, viabilizando sua execução
em dispositivos quânticos digitais.

\subsection{Implementação do algoritmo Variational Quantum Eigensolver}

A implementação do algoritmo VQE foi conduzida com base no fluxo
metodológico consolidado na literatura recente e no material
educacional disponibilizado pela IBM Quantum, sendo adaptada
para fins de consolidação conceitual e reprodutibilidade computacional.

\subsubsection{Construção do Hamiltoniano eletrônico}

A geometria molecular dos sistemas H$_2$ e LiH foi inicialmente
definida e utilizada para a construção do Hamiltoniano eletrônico
em segunda quantização, empregando o módulo Qiskit Nature.
Os integrais eletrônicos foram obtidos na base atômica mínima
STO-3G, escolha amplamente adotada em estudos introdutórios por
reduzir o custo computacional mantendo a estrutura física essencial
do problema.
O Hamiltoniano fermiónico resultante foi então mapeado para operadores
de spin por meio da transformação de Jordan–Wigner, permitindo sua
representação como combinação linear de operadores de Pauli
mensuráveis em dispositivos quânticos digitais.

\subsubsection{Definição do ansatz variacional}

Para a construção do estado variacional parametrizado, foi adotado
um \textit{hardware-efficient ansatz}, conforme proposto no tutorial
da IBM Quantum. Essa escolha justifica-se por sua adequação ao regime
NISQ, uma vez que tais circuitos apresentam baixa profundidade e
estrutura compatível com conectividades típicas de processadores
quânticos atuais.
O circuito parametrizado atua sobre o estado de referência de
Hartree–Fock, introduzindo rotações parametrizadas e camadas
de portas de entrelaçamento, produzindo estados candidatos ao
estado fundamental do sistema molecular. O número de camadas
foi mantido reduzido, de modo a equilibrar expressividade e
estabilidade numérica do processo de otimização.

\subsubsection{Loop híbrido quântico-clássico}

O núcleo do algoritmo consiste em um ciclo iterativo híbrido,
ilustrado na Figura~\ref{fig:vqe_cycle}, no qual:
\begin{itemize}
\item o circuito quântico prepara o estado variacional;
\item medições são realizadas para estimativa do valor esperado
da energia por meio da decomposição do Hamiltoniano em operadores
de Pauli;
\item um otimizador clássico atualiza os parâmetros variacionais;
\item o processo é repetido até a convergência da função custo.
\end{itemize}

Essa arquitetura híbrida explora a complementaridade entre a
capacidade de dispositivos quânticos prepararem estados altamente
entrelaçados e a eficiência de algoritmos clássicos na otimização
numérica, característica central dos algoritmos variacionais na
era NISQ.

\subsubsection{Otimização clássica}

A minimização da energia esperada foi realizada por meio do
otimizador clássico COBYLA (Constrained Optimization BY Linear
Approximations), conforme implementado no Qiskit.
A escolha do COBYLA justifica-se por tratar-se de um método
livre de derivadas, adequado para funções custo estimadas por
medições quânticas sujeitas a ruído estatístico. Esse método
apresenta boa estabilidade em problemas de baixa dimensionalidade,
como no caso das moléculas consideradas.
O critério de parada adotado baseou-se na convergência da energia
e no número máximo de iterações previamente definido.

\subsection{Execução das simulações}

As execuções foram realizadas utilizando o simulador de estado
vetorial do Qiskit Aer, permitindo a obtenção exata do valor
esperado da energia sem a introdução de ruído estatístico
associado a amostragem finita de medições.

Inicialmente, a metodologia foi reproduzida para a molécula
de hidrogênio (H$_2$), conforme apresentado no material de
referência da IBM Quantum, com o objetivo de validação da
implementação. Em seguida, o procedimento foi estendido para
o hidreto de lítio (LiH), permitindo avaliar o comportamento
do algoritmo diante de um sistema molecular de maior complexidade.
Os resultados foram analisados em termos de convergência do
processo variacional, estabilidade do otimizador e aproximação
obtida para a energia do estado fundamental, caracterizando a
implementação como uma prova de conceito do fluxo computacional
do VQE em ambiente controlado.